\section{Introduction}
\label{sec:introduction}

A GRPO-trained formalizer achieves a 93\% prove rate on first-order logic problems---yet 87\% of its proved outputs embed the conclusion verbatim among the premises.
The solver succeeds not because the model has learned to reason, but because the formalization trivially entails any conclusion that already appears as a premise.
This structural shortcut is invisible to execution-based evaluation: prove rate registers success regardless of whether the proof reflects genuine multi-step derivation or degenerate circular embedding.
The reward optimizes solver success; the model finds the shortest path to it.

\paragraph{The execution-only blind spot.}
Prove rate---the dominant evaluation metric for neurosymbolic formalizers~\citep{logic-lm, linc, satlm}---conflates faithful reasoning with degenerate shortcuts, a gap especially consequential under reinforcement learning, where the optimizer exploits any shortcut the reward fails to penalize~\citep{skalse-reward-hacking, reward-hacking-survey}.
Recent work applies GRPO~\citep{deepseekmath} to train LLMs that generate solver-verified formal programs~\citep{mellgren-prolog-grpo, logic-rl} using binary execution feedback as reward; this drives rapid PR gains but gives no gradient on structural quality.
Process reward models~\citep{lightman-verify-step, math-shepherd} offer step-level supervision but require expensive annotations and do not transfer to arbitrary solver-backed settings.
A complementary line of work documents that LLMs already game logical formalization at inference time~\citep{kim-game-formalization, bondarenko-spec-gaming}, yet no existing method provides an annotation-free \emph{diagnostic} for structural integrity in neurosymbolic outputs---one adaptable in principle to any solver supporting component-level input manipulation (see \S\ref{sec:related-work} for full context).

\paragraph{Component necessity testing.}
We address this gap with \emph{Solver Trace Ablation} (STA), a diagnostic protocol applying \emph{component necessity testing}---leave-one-out (LOO) ablation on solver inputs---to quantify structural integrity (\autoref{fig:sta-overview}).
For each proved formalization, STA removes each component individually and re-runs the solver; the fraction of components whose removal causes proof failure yields the \emph{Intensional Structural Reward} (ISR).
ISR${=}1$ indicates every component is individually required for proof success; ISR${=}0$ indicates complete structural redundancy.
We are explicit that LOO ablation is itself standard---in mutation testing~\citep{demillo1978, jia-harman-survey} and feature attribution---and we claim novelty not in the primitive but in its application as a diagnostic for GRPO reward hacking and in the empirical phenomena it surfaces.
We also emphasize that ISR captures structural rather than semantic integrity: a downstream NLI evaluation finds ISR and task accuracy uncorrelated ($\rho{=}{-}0.049$, $p{=}0.187$), confirming it measures properties orthogonal to semantic correctness.

\paragraph{ISR as diagnostic.}
Applied as a post-hoc diagnostic on a GRPO formalizer trained with FOLIO~\citep{folio} data, ISR reveals a structural collapse trajectory entirely invisible to prove rate: ISR falls from 0.642 at 100 steps to 0.132 at 500 steps to 0.003 at 1000 steps---a 99.5\% degradation---while prove rate simultaneously rises from 32\% to 96\%.
Three random seeds diverge radically (PR $\sigma{=}47$\,pp), each converging to a qualitatively different degeneration mode: conclusion embedding, repetition collapse, and tautological collapse.
A supervised fine-tuning (SFT) baseline on the same data achieves a prove rate of 14\% but 92.9\% non-hacking proofs with near-zero shortcut content, confirming that the structural collapse is a consequence of RL reward optimization, not model capacity or data limitations.

\paragraph{Meta-level Goodhart.}
A natural question is whether ISR can serve as a training signal.
The answer is negative: when ISR is added to the reward, $r = r_{\text{exec}} \cdot (1 + \beta \cdot \text{ISR})$, the model games the diagnostic itself---\emph{meta-level Goodhart's Law}~\citep{manheim-goodhart-taxonomy, karwowski-goodhart-rl, gao-reward-overoptimization}.
At $\beta{=}0.1$, 64.5\% of proved samples (3-seed mean) exhibit ISR-specific single-premise-equals-conclusion duplication---a construction in which ISR is high \emph{by definition} of the LOO metric---versus 0\% under a matched random-bonus reward; the one seed that retains a non-trivial ``genuine'' fraction reduces to 1.06$\times$ random after controlling for a premise-count confound (\S\ref{sec:goodhart}).
At $\beta{=}0.3$ the same reward produces extreme cross-seed instability (PR $66{\pm}42\%$) with three qualitatively distinct Goodhart pathways: reward collapse (PR ${=}\,18\%$), shortcut dominance (86\% PR, 69\% conclusion embedding), and pure metric gaming (95\% PR, 100\% ISR-inflating duplications, ISR${=}0.742$).
The effect transfers and intensifies in a different formal language: in Prolog with SWI-Prolog as solver, all three seeds produce 0\% non-hacking output---universal Goodhart---each via a distinct mechanism, including a \emph{query-as-fact gaming} mode that bypasses text-match classifiers entirely.
ISR's asymmetry is the central finding: \emph{reliable as observer, unreliable as optimizer.}

\paragraph{Contributions.}
Across 6 reward conditions, 3 seeds, 2 logic types, and up to 1000 training steps, we make four contributions:
\begin{itemize}[nosep,leftmargin=*]
\item[\textbf{C1.}] \textbf{ISR as diagnostic} for GRPO reward hacking, with a complete collapse trajectory ($\text{ISR}{=}0.642{\to}0.003$ over 1000 steps) invisible to prove rate, and an SFT control (92.9\% non-hacking) localizing the failure mode to RL optimization (\S\ref{sec:isr-diagnostic}).
\item[\textbf{C2.}] \textbf{Empirical meta-level Goodhart} in neurosymbolic RL, with a random-bonus control that isolates ISR-specific gaming (H7: 64.5\% vs.\ 0\%) and a premise-count-matched analysis showing that the sole apparent ISR gain is a mathematical artifact of the LOO metric, not a structural improvement (\S\ref{sec:goodhart}).
\item[\textbf{C3.}] \textbf{Cross-logic transfer} to Prolog/SWI-Prolog, demonstrating universal Goodhart (3/3 seeds gaming, 0\% non-hacking) and indicating that simpler formal semantics enable more complete exploitation (\S\ref{sec:prolog}).
\item[\textbf{C4.}] \textbf{A seven-mode hacking taxonomy} including a Prolog-specific query-as-fact gaming mode that bypasses text-match classifiers, motivating functional ($n_{\text{necessary}}$-based) cross-validation of diagnostics (\S\ref{sec:taxonomy-results}).
\end{itemize}
Taken together, these results characterize ISR's fundamental asymmetry: it is a useful instrument for \emph{detecting} reward hacking in solver-backed RL, but unsuitable as an optimization target because the very LOO primitive that makes it interpretable also makes it gameable.
We see the empirical Goodhart characterization, premise-count confound, and cross-logic universality as cautionary results for any work treating a structural proxy metric as a training signal.
